\chapter{Conclusions and Future Work}
\label{ch7}

\graphicspath{{Chapter8/}}


\GG{Software engineering education faces a structural tension that no single pedagogical intervention has fully resolved: the difficulty of teaching formal notation systems in ways that are both motivationally sustainable over the duration of a course and scalable enough to provide the personalized formative feedback that proficiency development requires. This thesis was designed to characterize, through empirical investigation, how two categories of technology-mediated intervention perform against that challenge in the specific domain of software modeling education. The gamification strand addressed the absence of rigorously evaluated tools for teaching BPMN and UML class diagrams through free-form, feedback-driven modeling environments — a gap that the literature mapping described in Section~\ref{sec:back_gamif_mde} made quantifiable: existing solutions were rare, relied predominantly on constrained element selection rather than free diagram construction, and had not been tested in longitudinal deployments that would allow the full intended mechanics to operate. The LLM strand addressed the corollary question of whether general-purpose language models could assist the automated evaluation and exercise-creation workflows that gamified tools require, at a moment when diagram generation had received far less empirical attention in the LLM-for-SE literature than code-centric tasks \cite{fan2023large, hou2023large}. Together, the eight studies reported here constitute a systematic multi-year attempt to characterize what is achievable within these constraints and what the field must still address.}

\GG{The four gamification studies, considered as a body of evidence rather than as individual experiments, support a coherent claim about when and how gamification produces benefits in modeling education. Across the preliminary BIPMIN study, the 2023 controlled BIPMIN experiment, the 2024 UMLegend experiment, and the 2023--2024 longitudinal BIPMIN deployment, the most consistently observed effect was an increase in the frequency with which students submitted their diagrams for automated evaluation: gamified students performed more checks per session than their non-gamified counterparts in every study that measured the metric, a behavioral pattern that directly reflects the mechanics design — if gamification feedback is triggered by evaluations, then more evaluations mean more exposure to corrective information. This elevated interaction with the feedback loop translated, in both controlled studies, into statistically significant improvements in diagram completeness and reductions in syntactic and semantic error counts \cite{garaccione_gamification_2024, garaccione_2026_gamification}, and extended in the longitudinal study to a statistically significant improvement in end-of-course exam and project grades relative to the preceding non-gamified cohort \cite{garaccione2024esem}. The evidence is therefore not that gamification is universally beneficial, but that gamification whose mechanics are tightly coupled to an automated evaluation engine — and which thereby converts every modeling decision into a feedback event — produces measurable improvements in both artifact quality and the development of modeling competence over time.}

\GG{What the gamification studies reveal about the design of effective mechanics is at least as important as the aggregate performance findings. Across all three tool experiments in which students rated individual mechanics, those providing direct element-level corrective guidance received the most consistently positive assessments: diagram coloring that highlights erroneous or close-matching elements, error lists that specify the type and location of each identified mistake, and completeness indicators that communicate progress distance rather than a binary pass-or-fail outcome. The mechanics designed for long-term engagement — level progression, experience points, avatar customization, and leaderboards — were appreciated considerably less across both single-session studies and, somewhat unexpectedly, in the longitudinal study as well, where leaderboard engagement fell short of expectations despite the semester-long deployment window. Interpreted through the lens of Self-Determination Theory, this asymmetry suggests that gamified modeling tools activate the competence need most reliably, since students can observe and correct specific skill gaps in real time, while the relatedness and autonomy needs that social and cosmetic mechanics are intended to serve require more carefully designed implementations than those developed here \cite{deci2000intrinsic, deci1985intrinsic}. The implication for future design is not that long-term mechanics are ineffective, but that their motivational contribution is more sensitive to implementation quality, social context, and deployment duration than the feedback mechanics that operate within a single exercise session, and that evaluating them meaningfully requires experimental designs that the present studies, including the longitudinal one, were not fully equipped to provide.}

\GG{The four LLM studies converge on an equally coherent, if more bounded, characterization of current model capabilities. The comparison between GPT-4 and a human solver on twenty class diagram exercises found that LLMs perform at a level statistically indistinguishable from a human on syntactic and pragmatic quality, while falling significantly behind on semantic error count and semantic distance from reference solutions \cite{debari2024evaluating}. The study on use case diagrams found no statistically significant differences in either completeness or redundancy rate, a result best understood as a consequence of the lower semantic complexity of use case notation: identifying actors and functional goals from natural language requires less domain inference than deriving association types, cardinality constraints, and attribute placement for class diagrams \cite{garaccione2025vega}. The multi-model comparative study, applying a role-based few-shot prompting strategy to GPT-4o, DeepSeek v3, Gemini 2.5 Pro, and Qwen~3 across fifteen exercises, confirmed that advanced prompting essentially eliminated syntax violations for three of the four models while leaving semantic difficulties — most notably incorrect association types and incorrect cardinality values — only partially resolved \cite{garaccione2025models}. Taken together, these findings point to a specific and consistent characterization of what general-purpose LLMs can and cannot do in this domain: they have internalized UML's structural grammar and can apply it reliably when prompted with appropriate formatting constraints, but they cannot yet reliably infer the domain-specific structural semantics — the relational logic that determines whether a relationship should be a composition or an aggregation, whether cardinality should be one-to-many or many-to-many — that constitutes the hardest part of class diagram generation and the part most consequential for educational use.}

\GG{The contribution of this thesis relative to prior work can be evaluated along three axes established by the background review. With respect to the literature on gamified tools for software modeling education (Section~\ref{sec:back_gamif_mde}), the thesis provides the first empirical demonstration that a free-form, gamified modeling environment produces significant improvements in diagram quality in both single-session and longitudinal deployments, and the first semester-length evidence that these benefits extend to end-of-course academic performance. No tool identified in the mapping review had been evaluated in a longitudinal controlled study, and none had been designed according to a framework that explicitly distinguishes exercise-specific from long-term mechanics — a distinction that the empirical results here validate as theoretically meaningful and practically important for experimental design. With respect to the automated evaluation literature (Section~\ref{sec:uml_eval}), the three-phase evaluation pipelines developed for BIPMIN and UMLegend extend the state of the art by integrating syntactic, semantic, and pragmatic assessment in a single workflow that is directly coupled to gamification feedback, supports multiple reference solutions and synonym-based matching, and generates per-element error messages derived from the quality taxonomy of Bolloju and Leung \cite{bolloju2006assisting}. This design provides more actionable formative information than the aggregate-score approaches of UMLGrade \cite{farias2020whats} or the scalar semantic-distance metric of Nikiforova et al.\ \cite{nikiforova2015approach}, and more flexibility than the single-reference, exact-name-matching pipelines found in the automated grading literature. With respect to the emerging literature on LLMs for software modeling tasks, the four studies in Chapter~6 provide the first systematic comparison of LLM-generated diagrams against student-produced ones across two notation types under controlled conditions, and the first multi-model evaluation applying a common advanced prompting strategy to class diagram generation, extending the single-model evaluations of prior work \cite{camara2023assessment, chen2023automated} to a comparative benchmark involving both proprietary and open-source architectures.}

\GG{The limitations of this research are most informatively discussed in terms of what they mean for the strength of the conclusions rather than as a defensive enumeration. The most consequential limitation — because it operates in the same direction as the positive findings — is the evaluation engine's dependence on teacher-defined synonym lists and structural alternatives. A semantically valid student response that falls outside the anticipated alternatives produces a false negative that simultaneously penalizes the student and inflates the measured error count, introducing noise into the completeness and error metrics that serve as the primary dependent variables in the gamification studies. The positive effects reported for these metrics are therefore lower bounds: a semantically flexible evaluation engine would likely yield larger effect sizes in the same experimental designs, because it would eliminate a source of noise that systematically attenuates the measured benefit of gamification. The persistent appearance of evaluation engine complaints as the most common improvement suggestion across open-ended questionnaire responses — in the 2023 BIPMIN study, the 2024 UMLegend study, and the longitudinal BIPMIN study alike — is the clearest evidence that this limitation was felt by students and affected their experience in ways the quantitative metrics alone do not fully capture. The external validity limitation is equally important to acknowledge: all gamification studies were conducted at Politecnico di Torino with Master's-level students who performed experiments at or near the end of their courses, when syntax competence had largely been established. The finding that gamification did not significantly reduce syntax errors in the 2023 controlled experiment, and showed only mixed effects in the longitudinal study, almost certainly reflects a ceiling effect in syntax knowledge rather than an intrinsic limitation of gamification as a pedagogical strategy, and studies conducted at the beginning of a modeling course would be needed to determine whether this pattern holds under conditions of genuine novice learning. The LLM studies are additionally constrained by small exercise samples of fifteen to twenty items per study, single human solvers providing the student performance baselines, and the use of interactive chat interfaces rather than controlled API pipelines; these constraints limit both statistical power and replicability, and the exercise sets drawn from Italian-language university curricula may not generalize to other linguistic or institutional contexts.}

\GG{For instructors and course designers, the most actionable finding from the gamification strand is the robustness of the feedback-mechanics effect: across every study in which individual mechanics were rated, error-level diagram coloring and element-specific error lists received the most consistently positive assessments, were associated with the largest measurable improvements in diagram quality, and were perceived as educationally useful regardless of whether deployment lasted a single session or an entire semester. This finding has a direct design implication: a gamified modeling tool that lacks these mechanics is unlikely to produce the quality improvements reported here, however elaborate its competitive or cosmetic features may be. The longitudinal study adds a necessary constraint for course designers: the statistically significant improvements in end-of-course grades described in Chapter~5 were observed only when the tool was available for the full semester, exercises were sequenced to introduce progressive challenge, and students had the opportunity to practice across a meaningful range of modeling scenarios \cite{garaccione2024esem}. Single-session deployments can improve within-session performance and provide valuable formative information, but they cannot replicate the grade-level outcomes that depend on sustained practice and skill accumulation over time. For tool developers, the TIGRE proof-of-concept suggests a practical and immediately usable workflow for the reference-solution creation bottleneck that currently constrains the scalability of automated evaluation: an LLM prompted with a role-based, structured instruction set and given two few-shot examples produces a first-pass solution in under a minute that contains most of the required classes and main associations, which an instructor then reviews and corrects using a visual editing interface, reducing the multi-hour effort currently required for manual reference solution creation to a lighter review-and-refinement task whose cognitive demand is concentrated on the semantic details — attribute types, cardinality values, association types — where current models still err systematically \cite{garaccione2026automated}.}

\GG{The future research directions motivated by this thesis form a coherent agenda rather than a set of independent suggestions. The most structurally pressing direction is the integration of the BIPMIN and UMLegend tools into a single multi-notation educational platform, a goal identified independently in the lessons learned sections of Chapters~3, 4, and 5 and explicitly requested by students in the longitudinal study's open-ended responses. The present situation — in which students interact with a gamified BPMN environment for process modeling activities and a separate gamified UML environment for class diagram activities, each with its own progression system, leaderboard, and avatar state — fragments the learning context in ways that undermine the coherence of the gamification experience and prevent the long-term mechanics from accumulating motivational weight across the modeling curriculum as a whole. Integration would also enable a more rigorous test of whether leaderboards and level systems produce meaningful engagement effects when they aggregate performance across a full course's worth of modeling activity, rather than within a single notation over a period where novelty effects confound the longitudinal signal. The evaluation engine limitation, which has been the single most persistent obstacle to the validity of the gamification findings, calls for an LLM-based semantic matching component that replaces normalized edit-distance comparison with contextual similarity assessment, directly addressing the synonym-completeness problem without requiring instructors to anticipate every valid label a student might use. The finding from Study~4 that DeepSeek, a model optimized for reasoning and code generation, produced the fewest semantic errors among the four evaluated architectures \cite{garaccione2025models} suggests that a purpose-specific or fine-tuned agent for association inference — the dimension on which all current general-purpose models fail most consistently — may be achievable with targeted training investment. The TIGRE assistant warrants full deployment and controlled evaluation: a study comparing the quality and instructor effort of reference solutions created with and without LLM assistance, across a representative sample of exercises and instructors, would establish whether the time-saving observed informally in the proof of concept holds under conditions that allow confounds to be isolated and would quantify the human correction effort required to bring LLM-generated solutions to the quality standard needed for reliable automated feedback. The avatar system requires dedicated design attention as an inclusivity problem: the concern raised by a student in the longitudinal study about stereotyping in algorithmically generated human avatars is not a minor usability defect but a design failure with equity implications, and future iterations should evaluate non-human or stylized representations that have been explicitly developed and assessed against inclusive design standards. Finally, and most directly, the observation that all gamification studies were conducted when students' syntax competence was already substantially established identifies the most educationally motivated replication: deploying BIPMIN or UMLegend at the beginning of a modeling course, before students have been exposed to formal notation instruction, would test whether gamified automated feedback can accelerate the formation of correct modeling habits and reduce the prior knowledge asymmetry that limits the conclusions of every study reported here.}

\GG{Several boundaries on the scope and generalizability of these findings must be stated at the outset. All empirical studies described in this thesis were conducted at Politecnico di Torino, in the context of Master's-level courses in management engineering and computer engineering. The student populations, exercise domains, and institutional structures reflect this specific context, and the results should not be assumed to transfer to undergraduate courses, professional or industrial training programs, or institutions with different curricular structures. The gamification studies evaluate BIPMIN and UMLegend in particular instantiations of the gamification framework; the framework itself is intended to generalize across modeling notations and educational contexts, but its transferability has not been empirically tested beyond the BPMN and UML class diagram domains covered here. The LLM studies are exploratory evaluations conducted with general-purpose models available during the 2023--2025 period; they characterize the capabilities and limitations of specific model families at a specific moment in a rapidly evolving technological landscape, and the findings about which error types are most prevalent and which prompting strategies are most effective should be interpreted as baseline measurements rather than as stable facts about LLM behavior. The TIGRE tool is evaluated through a two-exercise proof of concept rather than a controlled study, and its practical value for instructors requires further validation under more representative conditions.}

Software engineering education faces a structural tension that no single pedagogical intervention has fully resolved: the difficulty of teaching formal notation systems in ways that are both motivationally sustainable over the duration of a course and scalable enough to provide the personalized formative feedback that proficiency development requires. This thesis was designed to characterize, through empirical investigation, how two categories of technology-mediated intervention perform against that challenge in the specific domain of software modeling education. The gamification strand addressed the absence of rigorously evaluated tools for teaching BPMN and UML class diagrams through free-form, feedback-driven modeling environments — a gap that the literature mapping described in Section~\ref{sec:back_gamif_mde} made quantifiable: existing solutions were rare, relied predominantly on constrained element selection rather than free diagram construction, and had not been tested in longitudinal deployments that would allow the full intended mechanics to operate. The LLM strand addressed the corollary question of whether general-purpose language models could assist the automated evaluation and exercise-creation workflows that gamified tools require, at a moment when diagram generation had received far less empirical attention in the LLM-for-SE literature than code-centric tasks \cite{fan2023large, hou2023large}. Together, the eight studies reported here constitute a systematic multi-year attempt to characterize what is achievable within these constraints and what the field must still address.

The first research question asked whether gamification applied to software modeling exercises leads to measurable improvements in student performance, specifically in diagram productivity and quality, and the answer drawn from the controlled experiments in Chapters~3 and~4 is affirmative with an important qualification about the nature of the effect. Gamified students performed statistically significantly more automated evaluations of their diagrams than non-gamified counterparts in both controlled experiments, a behavioral change that directly reflects the coupling between the evaluation engine and the gamification mechanics: when each diagram submission triggers feedback that updates game state, students are incentivized to interact more with the evaluation loop, and more interactions mean more exposure to corrective information. This elevated engagement translated into statistically significant improvements in diagram completeness and reductions in both syntactic and semantic error counts in the 2024 UMLegend experiment \cite{garaccione_2026_gamification}, and into a positive effect on semantic correctness in the 2023 BIPMIN experiment \cite{garaccione_gamification_2024}. The qualification is equally important: gamification did not change the size of diagrams significantly, nor did it prompt students to introduce structural changes they would not otherwise have made. The mechanism is best characterized as refinement rather than expansion — students corrected and improved existing elements more diligently under gamified conditions, using the feedback loop as a guide for incremental improvement rather than as a signal to redesign diagrams wholesale. This distinction matters for interpreting what gamification actually does pedagogically: it sharpens attention to correctness within a student's current conceptual scope rather than broadening that scope.

The second research question asked whether sustained, longitudinal exposure to a gamified modeling tool can produce improvements in modeling competence that persist beyond any single session and manifest in end-of-course academic performance. The answer provided by Chapter~5 constitutes the strongest finding of the gamification strand: students in the 2023--2024 Information Systems course, who used the improved BIPMIN tool throughout the semester for progressively challenging BPMN exercises, achieved statistically significantly higher grades on both BPMN-related exam tasks and group project assignments than the preceding cohort who had learned the same material without the gamified tool \cite{garaccione2024esem}. The effect was most pronounced for the project component, where the magnitude of the improvement suggests that students who had practiced multi-entity, multi-pool processes over a sequence of five escalating exercises were substantially better equipped for the open-ended project requirements than their predecessors. The longitudinal analysis of exercise performance also illuminated the learning trajectory beneath this aggregate result: semantic error counts showed a statistically significant reduction from the first exercise to the last, indicating that repeated gamified practice helped students internalize domain representation conventions progressively over time. Completeness improvements over the same period were less consistent, however, with one exercise pair showing a statistically significant drop that likely reflects an ordering mismatch in the difficulty progression rather than a true regression in competence. The overall answer to the second question is therefore that longitudinal gamified practice can improve modeling competence and academic performance in measurable ways, but the design of the exercise sequence — specifically, the consistency of the difficulty gradient — materially affects whether the observed trajectory is one of steady growth or one punctuated by temporary setbacks.

The third research question asked how students perceive gamified modeling tools and which specific categories of mechanics most effectively support learning and engagement. The answer that emerges from post-study questionnaire analyses across all three tool studies is strikingly consistent and provides perhaps the most actionable design guidance in the thesis. Mechanics that provide direct element-level corrective guidance within an exercise session — diagram coloring that marks erroneous or close-matching elements, error lists that enumerate specific violations with descriptions, and completeness indicators that communicate remaining distance to exercise success — received the most positive TAM and GAMEX assessments across the 2023 BIPMIN study, the 2024 UMLegend study, and the longitudinal BIPMIN deployment, and their appreciation was robust regardless of whether the study lasted one session or an entire semester. The mechanics designed for long-term engagement — level progression, leaderboards, and avatar customization — were appreciated substantially less across all three experiments, including the longitudinal one where their benefits might have been expected to materialize given the semester-long window. A notable exception to this pattern is the improvement in GAMEX dimensions between the prototype study and the longitudinal deployment: enjoyment and creative thinking improved substantially with the redesigned tool, confirming that the framework changes described in Chapter~4 produced a genuine change in user experience even as the appreciation of competitive and cosmetic mechanics remained modest. Interpreted through Self-Determination Theory, the persistent asymmetry between feedback mechanics and long-term mechanics suggests that the competence need is the most reliably activated by the current designs, while the relatedness and autonomy needs targeted by social and cosmetic mechanics require implementations with stronger social infrastructure — cooperative challenges, visible peer comparison, or collective progression goals — than the present tools provided \cite{deci2000intrinsic, deci1985intrinsic}.

The fourth research question asked whether LLM agents can generate UML class diagrams and use case diagrams whose quality is comparable to that of student-produced artifacts, and the answer differs significantly by notation type. For class diagrams, the comparison between GPT-4 and a human solver on twenty exercises found a statistically significant inferiority of the LLM on semantic error count and semantic distance from expert reference solutions, while differences on syntactic and pragmatic quality were not statistically significant \cite{debari2024evaluating}. The model produced diagrams that were structurally well-formed and internally consistent, but systematically failed to capture domain-specific constraints — correct association types, cardinality values, and the distinction between concepts that should be classes versus attributes — at the same rate as a trained student. For use case diagrams, in contrast, no statistically significant differences were found between the LLM and the human solver on either completeness rate or redundancy rate \cite{garaccione2025vega}, a result that reflects the structurally simpler demands of the notation: identifying actors and user goals from natural language is a task LLMs perform at a student-comparable level, while the deeper domain-inferential reasoning required for class diagrams is not. The answer to the fourth research question is therefore notation-dependent: LLMs are viable generators of use case diagram drafts of student-comparable quality, but their class diagram output remains significantly weaker than student-produced work on the semantically demanding dimensions, making it more suitable as a starting point for human refinement than as a direct substitute for trained student or instructor-validated solutions.

The fifth research question asked whether an LLM-based tool can reduce the effort required by educators to create exercise reference solutions, the bottleneck that currently limits the scalability of automated formative feedback. The proof-of-concept evaluation of TIGRE in Chapter~6 provides an encouraging initial answer: the tool generated first-pass solutions for two past exam exercises in under a minute of processing time, identifying most of the required classes and a portion of the expected associations, and produced output directly importable into the Apollon visual editing environment for human review and correction \cite{garaccione2026automated}. Class identification was relatively strong — seven of ten and four of seven classes were judged valid matches in the two exercises respectively — while attribute completeness and association correctness were weaker, with cardinality values in particular being systematically incorrect across both generated diagrams. The practical implication is that LLM assistance concentrates the remaining instructor effort on the semantic details where models consistently err, rather than on the structural scaffolding, which the model handles acceptably. Whether this redistribution of effort constitutes a net reduction in total instructor time compared to manual creation from scratch cannot be established from a two-exercise proof of concept alone; a controlled comparative study with a representative sample of instructors and exercises is the necessary next step. On current evidence, however, the answer to the fifth research question is that LLM-based tool support can generate a usable structural scaffold for reference solutions quickly enough to make a review-and-correct workflow a credible alternative to full manual creation, provided instructors attend carefully to the association-related elements that current models routinely mishandle.

The sixth research question asked how different LLM architectures compare in diagram generation quality and whether role-based few-shot prompting improves output quality beyond zero-shot approaches. The multi-model comparative study in Chapter~6 provides a differentiated answer across three quality dimensions \cite{garaccione2025models}. For syntax quality, ChatGPT and Gemini produced diagrams with essentially zero violations under the structured few-shot prompt, DeepSeek showed minimal infractions, and Qwen stood as a clear outlier — its systematic use of collection types and foreign-key attribute patterns to represent associations violated UML structural conventions explicitly excluded in the prompt, suggesting that its training for this task type is less mature than the other three models. For semantic quality, the ordering inverts in a diagnostically informative way: DeepSeek, optimized for code and structured reasoning, produced the fewest semantic errors on average, indicating that its inferential strengths translate to an advantage when deriving implicit domain constraints from natural language descriptions. Diagram completeness was broadly comparable across all four models, with proprietary models performing marginally but not significantly better than open-source alternatives. On the prompting strategy question, the most direct evidence for the benefit of advanced prompting comes from comparing ChatGPT's performance between Study~1 and Study~4: a shift from a zero-shot single-sentence prompt to a role-based few-shot prompt reduced average syntax errors from approximately 0.9 to 0.0 for the same model family, demonstrating that prompt engineering can eliminate the syntax gap that zero-shot approaches leave open, even if it cannot close the semantic gap that persists under all current strategies.

The contribution of this thesis relative to prior work can be evaluated along three axes established by the background review. With respect to the literature on gamified tools for software modeling education (Section~\ref{sec:back_gamif_mde}), the thesis provides the first empirical demonstration that a free-form, gamified modeling environment produces significant improvements in diagram quality in both single-session and longitudinal deployments, and the first semester-length evidence that these benefits extend to end-of-course academic performance. No tool identified in the mapping review had been evaluated in a longitudinal controlled study, and none had been designed according to a framework that explicitly distinguishes exercise-specific from long-term mechanics — a distinction that the empirical results here validate as theoretically meaningful and practically important for experimental design. With respect to the automated evaluation literature (Section~\ref{sec:uml_eval}), the three-phase evaluation pipelines developed for BIPMIN and UMLegend extend the state of the art by integrating syntactic, semantic, and pragmatic assessment in a single workflow that is directly coupled to gamification feedback, supports multiple reference solutions and synonym-based matching, and generates per-element error messages derived from the quality taxonomy of Bolloju and Leung \cite{bolloju2006assisting}. This design provides more actionable formative information than the aggregate-score approaches of UMLGrade \cite{farias2020whats} or the scalar semantic-distance metric of Nikiforova et al.\ \cite{nikiforova2015approach}, and more flexibility than the single-reference, exact-name-matching pipelines found in the automated grading literature. With respect to the emerging literature on LLMs for software modeling tasks, the four studies in Chapter~6 provide the first systematic comparison of LLM-generated diagrams against student-produced ones across two notation types under controlled conditions, and the first multi-model evaluation applying a common advanced prompting strategy to class diagram generation, extending the single-model evaluations of prior work \cite{camara2023assessment, chen2023automated} to a comparative benchmark involving both proprietary and open-source architectures.

The limitations of this research are most informatively discussed in terms of what they mean for the strength of the conclusions rather than as a defensive enumeration. The most consequential limitation — because it operates in the same direction as the positive findings — is the evaluation engine's dependence on teacher-defined synonym lists and structural alternatives. A semantically valid student response that falls outside the anticipated alternatives produces a false negative that simultaneously penalizes the student and inflates the measured error count, introducing noise into the completeness and error metrics that are the primary dependent variables in the gamification studies. The positive effects reported for these metrics are therefore lower bounds on what a semantically flexible evaluation engine might produce: eliminating this source of noise would be expected to increase rather than reverse the observed effect sizes. The persistent appearance of evaluation engine complaints as the most common improvement suggestion across open-ended questionnaire responses — in the 2023 BIPMIN study, the 2024 UMLegend study, and the longitudinal BIPMIN study — confirms that this limitation affected student experience in ways the quantitative metrics alone do not capture. The external validity limitation is equally important: all gamification studies were conducted at Politecnico di Torino with Master's-level students performing experiments at or near the end of their courses, when syntax competence was already substantially established. The near-null effect of gamification on syntactic error reduction almost certainly reflects a ceiling effect on syntax knowledge rather than an intrinsic limitation of the approach, and the absence of studies conducted at the beginning of a modeling course means the most educationally motivated use case remains untested. The LLM studies are further constrained by small exercise samples of fifteen to twenty items per study, single human solvers providing the student baselines, and the use of interactive chat interfaces rather than controlled API pipelines; these constraints limit both statistical power and replicability, and the Italian-language exercise sets may not generalize to other linguistic or institutional contexts.

For instructors and course designers, the most actionable finding from the gamification strand is the robustness of the feedback-mechanics effect: across every study in which individual mechanics were rated, error-level diagram coloring and element-specific error lists were the most consistently appreciated, associated with the largest improvements in diagram quality, and perceived as educationally useful regardless of deployment duration. A gamified modeling tool that lacks these mechanics is unlikely to produce the quality improvements reported here, however elaborate its competitive or cosmetic features may be. The longitudinal study adds a necessary constraint for course designers: the improvements in academic performance observed in Chapter~5 depended on full-semester availability, progressive exercise difficulty, and student exposure to a meaningful range of modeling scenarios \cite{garaccione2024esem}. Single-session deployments can improve within-session performance and generate useful formative data, but they cannot replicate the grade-level outcomes that depend on sustained and cumulative practice. For tool developers, the TIGRE proof-of-concept suggests a credible workflow for the reference-solution creation bottleneck: a role-based few-shot prompt produces a structurally usable first-pass solution in under a minute, which an instructor then refines through a visual editing interface, concentrating human effort on the association semantics that current models handle poorly while delegating the more straightforward class and attribute identification to the model \cite{garaccione2026automated}. For educators considering LLM integration in modeling courses, the Study~4 finding that model selection should depend on the primary quality criterion — proprietary models for syntax reliability and broad coverage, reasoning-optimized open-source models for semantic precision on smaller diagrams — provides a first empirically grounded basis for practical tool choices.

The future research directions motivated by this thesis form a coherent agenda grounded in the specific open questions the studies leave behind. The most structurally pressing direction is the integration of the BIPMIN and UMLegend tools into a single multi-notation educational platform, a goal identified independently in the lessons learned sections of Chapters~3, 4, and~5 and explicitly requested by students in the longitudinal study's open-ended responses. The present situation — two separate gamified environments for two modeling notations, each with its own progression system, leaderboard, and avatar state — fragments the learning context in ways that prevent the long-term mechanics from accumulating motivational weight across the full modeling curriculum. Integration would also enable a controlled replication of the third research question under richer conditions: with mechanics aggregating engagement across notations, the underperformance of leaderboards and level systems observed in single-notation deployments could be re-evaluated under conditions that more closely match the system's intended design. The evaluation engine limitation calls for an LLM-based semantic matching component to replace normalized edit-distance comparison, directly addressing the synonym-completeness problem that introduced noise into all four gamification studies; the Study~4 evidence that DeepSeek's reasoning-oriented training produces superior semantic error rates \cite{garaccione2025models} suggests that association inference specifically, the dimension where all general-purpose models fail most consistently, may be tractable with a purpose-fine-tuned agent. The TIGRE assistant requires a controlled comparative study comparing reference-solution quality and instructor effort with and without LLM assistance, across a representative exercise sample, to determine whether the structural benefit observed in the proof of concept holds at scale. The avatar system requires dedicated redesign as an inclusivity problem: the concern raised in the longitudinal study about stereotyping in algorithmically generated avatars is a design failure with equity implications that stylized or non-human avatar representations could address. Finally, the fourth research question about LLM diagram quality was answered only at the end of the curriculum, when student baselines reflect substantial prior training; a study comparing LLM-generated and student-produced diagrams at the beginning of a modeling course — where both baselines are lower and the comparison is more educationally relevant — would test whether LLMs can serve as scaffolding for novice learners in ways that the present studies, which used advanced students as the comparison point, cannot assess.